% =====================================================================
% PERTENECE AL CAPITULO 3 (metodologia.tex), NO al 4.
% Estaba al final de resultados_notas_ampliacion.tex y se separo el
% 2026-09-28, al incorporar ese archivo a resultados.tex: de haber quedado
% adentro, esta subseccion habria aparecido dentro del capitulo de
% resultados.
%
% NO esta incluida en ningun \input todavia. Para incorporarla, agregar en
% metodologia.tex, donde corresponda:
%     % =====================================================================
% PERTENECE AL CAPITULO 3 (metodologia.tex), NO al 4.
% Estaba al final de resultados_notas_ampliacion.tex y se separo el
% 2026-09-28, al incorporar ese archivo a resultados.tex: de haber quedado
% adentro, esta subseccion habria aparecido dentro del capitulo de
% resultados.
%
% NO esta incluida en ningun \input todavia. Para incorporarla, agregar en
% metodologia.tex, donde corresponda:
%     % =====================================================================
% PERTENECE AL CAPITULO 3 (metodologia.tex), NO al 4.
% Estaba al final de resultados_notas_ampliacion.tex y se separo el
% 2026-09-28, al incorporar ese archivo a resultados.tex: de haber quedado
% adentro, esta subseccion habria aparecido dentro del capitulo de
% resultados.
%
% NO esta incluida en ningun \input todavia. Para incorporarla, agregar en
% metodologia.tex, donde corresponda:
%     % =====================================================================
% PERTENECE AL CAPITULO 3 (metodologia.tex), NO al 4.
% Estaba al final de resultados_notas_ampliacion.tex y se separo el
% 2026-09-28, al incorporar ese archivo a resultados.tex: de haber quedado
% adentro, esta subseccion habria aparecido dentro del capitulo de
% resultados.
%
% NO esta incluida en ningun \input todavia. Para incorporarla, agregar en
% metodologia.tex, donde corresponda:
%     \input{metodologia_corpus_notas}
% =====================================================================

\subsection{Construcción del corpus de notas: criterios de admisión}
\label{subsec:met_admision}

La validez de todo el frente de notas depende de que cada nota esté correctamente atribuida a su familia y provenga de una fuente verificable. Durante el trabajo se detectaron errores en ambos aspectos, y las reglas que se describen aquí se adoptaron como consecuencia de esos hallazgos, no de forma anticipada.

\subsubsection*{Qué cuenta como un texto distinto}

Dos notas se consideran la misma \textbf{plantilla} cuando su similitud de coseno sobre TF-IDF de $n$-gramas de caracteres (3 a 5) supera 0,90, y las plantillas se forman como componentes conexas de esa relación. El umbral no se eligió por conveniencia: es el mismo que separa los grupos que el protocolo P2 no reparte entre entrenamiento y prueba, de modo que \textbf{la unidad con que se cuenta el corpus y la unidad con que se evalúa son la misma}. Toda cifra de tamaño del corpus se reporta por lo tanto en dos magnitudes, notas y plantillas, porque la segunda es la que gobierna el rendimiento.

El criterio tiene un caso de frontera que conviene declarar. WASTEDLOCKER reúne cuatro notas que, leídas, son evidentemente el mismo molde: cambian la víctima y las direcciones de contacto. Sus similitudes internas, sin embargo, son 0,886, 0,888, 0,896, 0,898, 0,900 y 0,974, de modo que el criterio automático las cuenta como \textbf{tres} plantillas y no como una. La familia ilustra que el umbral, siendo necesario para operar, no coincide siempre con el juicio de un lector, y que el conteo de plantillas debe entenderse como una definición operativa y no como una verdad sobre el malware.

\subsubsection*{Verificación de cada nota candidata}

Toda nota candidata a incorporarse se somete al mismo criterio antes de contarse: si al añadirla el número de plantillas no aumenta, la nota es una variante trivial de material ya presente y \textbf{no se incorpora}. Esta regla es consecuencia directa del resultado de la curva de aprendizaje, según la cual lo que mueve el rendimiento son los textos distintos y no el volumen. En la práctica descartó una proporción alta de las candidatas reunidas: repositorios distintos publican con frecuencia exactamente el mismo texto.

Esa redundancia entre fuentes, que es un obstáculo para ampliar el corpus, resultó ser una ventaja para validarlo. Al contrastar el corpus contra catálogos independientes se encontró que varias de sus notas aparecen publicadas con similitudes de 0,979 a 1,000 en fuentes que no participaron de su recolección, lo que \textbf{corrobora que los textos son auténticos} y no reconstrucciones del recolector.

\subsubsection*{Homónimos y la distinción entre campaña y familia}

El error más costoso detectado no fue de medición sino de etiquetado, y se repitió en tres formas distintas:

\begin{itemize}
  \item \textbf{Homónimo de familia.} Dos notas atribuidas a MEDUZALOCKER pertenecían en realidad a \textit{Medusa}, una familia distinta cuyo nombre se parece. Otra, atribuida a CRYPTOLOCKER, pertenecía a \textit{Crypt0l0cker}, que es un alias de TorrentLocker.
  \item \textbf{Contaminación entre linajes.} Dos notas atribuidas a HELLOKITTY compartían el párrafo de cierre con doce notas de DHARMA y ninguna de su familia declarada. Su retiro eliminó 146 confusiones desde DHARMA y 20 desde PHOBOS en la matriz de confusión.
  \item \textbf{Campaña frente a familia.} Los sucesores de WASTEDLOCKER atribuidos al mismo actor son \textit{familias distintas}, y sus notas no pertenecen al corpus de WASTEDLOCKER aunque la atribución del actor sea correcta.
\end{itemize}

De ahí la regla adoptada: \textbf{el corpus se etiqueta por familia —entendida como el mismo linaje de código y de plantilla— y no por actor ni por campaña}, y ante nombres parecidos se exige una fuente que distinga explícitamente entre ambos. El emparejamiento automático de nombres se evitó en favor de una tabla de equivalencias escrita a mano, porque un normalizador de cadenas es exactamente el mecanismo por el que se cuelan estos tres errores.

\subsubsection*{Fuentes y trazabilidad}

Cada nota lleva registrada su procedencia en un manifiesto: la fuente, la URL cuando existe, el nombre original del archivo cuando se conserva y la forma en que se obtuvo el texto. Las fuentes se admiten con criterios explícitos —autoría verificable, dominio no susceptible de caducar y ser recomprado, preferencia por el proveedor original sobre agregadores comerciales— y se excluyen los repositorios de muestras ejecutables. Al cierre del corpus de 149 notas, \textbf{la totalidad tiene fuente verificable}, frente a las 47 sin procedencia registrada que había en versiones anteriores.

Cuando una nota se transcribe de una imagen, el texto se registra literalmente, incluidas las erratas de la fuente y las truncaduras que ella misma introduce, porque corregirlas produciría un texto que nunca existió. Los identificadores criptográficos transcritos así se verifican además \textbf{mediante su propia suma de comprobación}: una dirección de Bitcoin o de Monero incorpora dígitos de control, de modo que un solo carácter mal leído invalida la verificación. El procedimiento detectó dos caracteres incorrectos en una transcripción externa que a simple vista parecía correcta.

\subsubsection*{Limitación que se declara}

Ninguno de estos controles convierte al corpus en una muestra representativa. Las notas provienen de repositorios públicos, que sobre-representan familias analizadas por la industria y sub-representan campañas recientes o de bajo perfil; el propio proceso de catalogación elimina metadatos —el nombre original del archivo, en particular— que en un incidente real sí están disponibles. Los resultados de este trabajo deben leerse sobre esa base: \textbf{un corpus público, auditado y trazable, pero no una muestra aleatoria del fenómeno}.

% =====================================================================

\subsection{Construcción del corpus de notas: criterios de admisión}
\label{subsec:met_admision}

La validez de todo el frente de notas depende de que cada nota esté correctamente atribuida a su familia y provenga de una fuente verificable. Durante el trabajo se detectaron errores en ambos aspectos, y las reglas que se describen aquí se adoptaron como consecuencia de esos hallazgos, no de forma anticipada.

\subsubsection*{Qué cuenta como un texto distinto}

Dos notas se consideran la misma \textbf{plantilla} cuando su similitud de coseno sobre TF-IDF de $n$-gramas de caracteres (3 a 5) supera 0,90, y las plantillas se forman como componentes conexas de esa relación. El umbral no se eligió por conveniencia: es el mismo que separa los grupos que el protocolo P2 no reparte entre entrenamiento y prueba, de modo que \textbf{la unidad con que se cuenta el corpus y la unidad con que se evalúa son la misma}. Toda cifra de tamaño del corpus se reporta por lo tanto en dos magnitudes, notas y plantillas, porque la segunda es la que gobierna el rendimiento.

El criterio tiene un caso de frontera que conviene declarar. WASTEDLOCKER reúne cuatro notas que, leídas, son evidentemente el mismo molde: cambian la víctima y las direcciones de contacto. Sus similitudes internas, sin embargo, son 0,886, 0,888, 0,896, 0,898, 0,900 y 0,974, de modo que el criterio automático las cuenta como \textbf{tres} plantillas y no como una. La familia ilustra que el umbral, siendo necesario para operar, no coincide siempre con el juicio de un lector, y que el conteo de plantillas debe entenderse como una definición operativa y no como una verdad sobre el malware.

\subsubsection*{Verificación de cada nota candidata}

Toda nota candidata a incorporarse se somete al mismo criterio antes de contarse: si al añadirla el número de plantillas no aumenta, la nota es una variante trivial de material ya presente y \textbf{no se incorpora}. Esta regla es consecuencia directa del resultado de la curva de aprendizaje, según la cual lo que mueve el rendimiento son los textos distintos y no el volumen. En la práctica descartó una proporción alta de las candidatas reunidas: repositorios distintos publican con frecuencia exactamente el mismo texto.

Esa redundancia entre fuentes, que es un obstáculo para ampliar el corpus, resultó ser una ventaja para validarlo. Al contrastar el corpus contra catálogos independientes se encontró que varias de sus notas aparecen publicadas con similitudes de 0,979 a 1,000 en fuentes que no participaron de su recolección, lo que \textbf{corrobora que los textos son auténticos} y no reconstrucciones del recolector.

\subsubsection*{Homónimos y la distinción entre campaña y familia}

El error más costoso detectado no fue de medición sino de etiquetado, y se repitió en tres formas distintas:

\begin{itemize}
  \item \textbf{Homónimo de familia.} Dos notas atribuidas a MEDUZALOCKER pertenecían en realidad a \textit{Medusa}, una familia distinta cuyo nombre se parece. Otra, atribuida a CRYPTOLOCKER, pertenecía a \textit{Crypt0l0cker}, que es un alias de TorrentLocker.
  \item \textbf{Contaminación entre linajes.} Dos notas atribuidas a HELLOKITTY compartían el párrafo de cierre con doce notas de DHARMA y ninguna de su familia declarada. Su retiro eliminó 146 confusiones desde DHARMA y 20 desde PHOBOS en la matriz de confusión.
  \item \textbf{Campaña frente a familia.} Los sucesores de WASTEDLOCKER atribuidos al mismo actor son \textit{familias distintas}, y sus notas no pertenecen al corpus de WASTEDLOCKER aunque la atribución del actor sea correcta.
\end{itemize}

De ahí la regla adoptada: \textbf{el corpus se etiqueta por familia —entendida como el mismo linaje de código y de plantilla— y no por actor ni por campaña}, y ante nombres parecidos se exige una fuente que distinga explícitamente entre ambos. El emparejamiento automático de nombres se evitó en favor de una tabla de equivalencias escrita a mano, porque un normalizador de cadenas es exactamente el mecanismo por el que se cuelan estos tres errores.

\subsubsection*{Fuentes y trazabilidad}

Cada nota lleva registrada su procedencia en un manifiesto: la fuente, la URL cuando existe, el nombre original del archivo cuando se conserva y la forma en que se obtuvo el texto. Las fuentes se admiten con criterios explícitos —autoría verificable, dominio no susceptible de caducar y ser recomprado, preferencia por el proveedor original sobre agregadores comerciales— y se excluyen los repositorios de muestras ejecutables. Al cierre del corpus de 149 notas, \textbf{la totalidad tiene fuente verificable}, frente a las 47 sin procedencia registrada que había en versiones anteriores.

Cuando una nota se transcribe de una imagen, el texto se registra literalmente, incluidas las erratas de la fuente y las truncaduras que ella misma introduce, porque corregirlas produciría un texto que nunca existió. Los identificadores criptográficos transcritos así se verifican además \textbf{mediante su propia suma de comprobación}: una dirección de Bitcoin o de Monero incorpora dígitos de control, de modo que un solo carácter mal leído invalida la verificación. El procedimiento detectó dos caracteres incorrectos en una transcripción externa que a simple vista parecía correcta.

\subsubsection*{Limitación que se declara}

Ninguno de estos controles convierte al corpus en una muestra representativa. Las notas provienen de repositorios públicos, que sobre-representan familias analizadas por la industria y sub-representan campañas recientes o de bajo perfil; el propio proceso de catalogación elimina metadatos —el nombre original del archivo, en particular— que en un incidente real sí están disponibles. Los resultados de este trabajo deben leerse sobre esa base: \textbf{un corpus público, auditado y trazable, pero no una muestra aleatoria del fenómeno}.

% =====================================================================

\subsection{Construcción del corpus de notas: criterios de admisión}
\label{subsec:met_admision}

La validez de todo el frente de notas depende de que cada nota esté correctamente atribuida a su familia y provenga de una fuente verificable. Durante el trabajo se detectaron errores en ambos aspectos, y las reglas que se describen aquí se adoptaron como consecuencia de esos hallazgos, no de forma anticipada.

\subsubsection*{Qué cuenta como un texto distinto}

Dos notas se consideran la misma \textbf{plantilla} cuando su similitud de coseno sobre TF-IDF de $n$-gramas de caracteres (3 a 5) supera 0,90, y las plantillas se forman como componentes conexas de esa relación. El umbral no se eligió por conveniencia: es el mismo que separa los grupos que el protocolo P2 no reparte entre entrenamiento y prueba, de modo que \textbf{la unidad con que se cuenta el corpus y la unidad con que se evalúa son la misma}. Toda cifra de tamaño del corpus se reporta por lo tanto en dos magnitudes, notas y plantillas, porque la segunda es la que gobierna el rendimiento.

El criterio tiene un caso de frontera que conviene declarar. WASTEDLOCKER reúne cuatro notas que, leídas, son evidentemente el mismo molde: cambian la víctima y las direcciones de contacto. Sus similitudes internas, sin embargo, son 0,886, 0,888, 0,896, 0,898, 0,900 y 0,974, de modo que el criterio automático las cuenta como \textbf{tres} plantillas y no como una. La familia ilustra que el umbral, siendo necesario para operar, no coincide siempre con el juicio de un lector, y que el conteo de plantillas debe entenderse como una definición operativa y no como una verdad sobre el malware.

\subsubsection*{Verificación de cada nota candidata}

Toda nota candidata a incorporarse se somete al mismo criterio antes de contarse: si al añadirla el número de plantillas no aumenta, la nota es una variante trivial de material ya presente y \textbf{no se incorpora}. Esta regla es consecuencia directa del resultado de la curva de aprendizaje, según la cual lo que mueve el rendimiento son los textos distintos y no el volumen. En la práctica descartó una proporción alta de las candidatas reunidas: repositorios distintos publican con frecuencia exactamente el mismo texto.

Esa redundancia entre fuentes, que es un obstáculo para ampliar el corpus, resultó ser una ventaja para validarlo. Al contrastar el corpus contra catálogos independientes se encontró que varias de sus notas aparecen publicadas con similitudes de 0,979 a 1,000 en fuentes que no participaron de su recolección, lo que \textbf{corrobora que los textos son auténticos} y no reconstrucciones del recolector.

\subsubsection*{Homónimos y la distinción entre campaña y familia}

El error más costoso detectado no fue de medición sino de etiquetado, y se repitió en tres formas distintas:

\begin{itemize}
  \item \textbf{Homónimo de familia.} Dos notas atribuidas a MEDUZALOCKER pertenecían en realidad a \textit{Medusa}, una familia distinta cuyo nombre se parece. Otra, atribuida a CRYPTOLOCKER, pertenecía a \textit{Crypt0l0cker}, que es un alias de TorrentLocker.
  \item \textbf{Contaminación entre linajes.} Dos notas atribuidas a HELLOKITTY compartían el párrafo de cierre con doce notas de DHARMA y ninguna de su familia declarada. Su retiro eliminó 146 confusiones desde DHARMA y 20 desde PHOBOS en la matriz de confusión.
  \item \textbf{Campaña frente a familia.} Los sucesores de WASTEDLOCKER atribuidos al mismo actor son \textit{familias distintas}, y sus notas no pertenecen al corpus de WASTEDLOCKER aunque la atribución del actor sea correcta.
\end{itemize}

De ahí la regla adoptada: \textbf{el corpus se etiqueta por familia —entendida como el mismo linaje de código y de plantilla— y no por actor ni por campaña}, y ante nombres parecidos se exige una fuente que distinga explícitamente entre ambos. El emparejamiento automático de nombres se evitó en favor de una tabla de equivalencias escrita a mano, porque un normalizador de cadenas es exactamente el mecanismo por el que se cuelan estos tres errores.

\subsubsection*{Fuentes y trazabilidad}

Cada nota lleva registrada su procedencia en un manifiesto: la fuente, la URL cuando existe, el nombre original del archivo cuando se conserva y la forma en que se obtuvo el texto. Las fuentes se admiten con criterios explícitos —autoría verificable, dominio no susceptible de caducar y ser recomprado, preferencia por el proveedor original sobre agregadores comerciales— y se excluyen los repositorios de muestras ejecutables. Al cierre del corpus de 149 notas, \textbf{la totalidad tiene fuente verificable}, frente a las 47 sin procedencia registrada que había en versiones anteriores.

Cuando una nota se transcribe de una imagen, el texto se registra literalmente, incluidas las erratas de la fuente y las truncaduras que ella misma introduce, porque corregirlas produciría un texto que nunca existió. Los identificadores criptográficos transcritos así se verifican además \textbf{mediante su propia suma de comprobación}: una dirección de Bitcoin o de Monero incorpora dígitos de control, de modo que un solo carácter mal leído invalida la verificación. El procedimiento detectó dos caracteres incorrectos en una transcripción externa que a simple vista parecía correcta.

\subsubsection*{Limitación que se declara}

Ninguno de estos controles convierte al corpus en una muestra representativa. Las notas provienen de repositorios públicos, que sobre-representan familias analizadas por la industria y sub-representan campañas recientes o de bajo perfil; el propio proceso de catalogación elimina metadatos —el nombre original del archivo, en particular— que en un incidente real sí están disponibles. Los resultados de este trabajo deben leerse sobre esa base: \textbf{un corpus público, auditado y trazable, pero no una muestra aleatoria del fenómeno}.

% =====================================================================

\subsection{Construcción del corpus de notas: criterios de admisión}
\label{subsec:met_admision}

La validez de todo el frente de notas depende de que cada nota esté correctamente atribuida a su familia y provenga de una fuente verificable. Durante el trabajo se detectaron errores en ambos aspectos, y las reglas que se describen aquí se adoptaron como consecuencia de esos hallazgos, no de forma anticipada.

\subsubsection*{Qué cuenta como un texto distinto}

Dos notas se consideran la misma \textbf{plantilla} cuando su similitud de coseno sobre TF-IDF de $n$-gramas de caracteres (3 a 5) supera 0,90, y las plantillas se forman como componentes conexas de esa relación. El umbral no se eligió por conveniencia: es el mismo que separa los grupos que el protocolo P2 no reparte entre entrenamiento y prueba, de modo que \textbf{la unidad con que se cuenta el corpus y la unidad con que se evalúa son la misma}. Toda cifra de tamaño del corpus se reporta por lo tanto en dos magnitudes, notas y plantillas, porque la segunda es la que gobierna el rendimiento.

El criterio tiene un caso de frontera que conviene declarar. WASTEDLOCKER reúne cuatro notas que, leídas, son evidentemente el mismo molde: cambian la víctima y las direcciones de contacto. Sus similitudes internas, sin embargo, son 0,886, 0,888, 0,896, 0,898, 0,900 y 0,974, de modo que el criterio automático las cuenta como \textbf{tres} plantillas y no como una. La familia ilustra que el umbral, siendo necesario para operar, no coincide siempre con el juicio de un lector, y que el conteo de plantillas debe entenderse como una definición operativa y no como una verdad sobre el malware.

\subsubsection*{Verificación de cada nota candidata}

Toda nota candidata a incorporarse se somete al mismo criterio antes de contarse: si al añadirla el número de plantillas no aumenta, la nota es una variante trivial de material ya presente y \textbf{no se incorpora}. Esta regla es consecuencia directa del resultado de la curva de aprendizaje, según la cual lo que mueve el rendimiento son los textos distintos y no el volumen. En la práctica descartó una proporción alta de las candidatas reunidas: repositorios distintos publican con frecuencia exactamente el mismo texto.

Esa redundancia entre fuentes, que es un obstáculo para ampliar el corpus, resultó ser una ventaja para validarlo. Al contrastar el corpus contra catálogos independientes se encontró que varias de sus notas aparecen publicadas con similitudes de 0,979 a 1,000 en fuentes que no participaron de su recolección, lo que \textbf{corrobora que los textos son auténticos} y no reconstrucciones del recolector.

\subsubsection*{Homónimos y la distinción entre campaña y familia}

El error más costoso detectado no fue de medición sino de etiquetado, y se repitió en tres formas distintas:

\begin{itemize}
  \item \textbf{Homónimo de familia.} Dos notas atribuidas a MEDUZALOCKER pertenecían en realidad a \textit{Medusa}, una familia distinta cuyo nombre se parece. Otra, atribuida a CRYPTOLOCKER, pertenecía a \textit{Crypt0l0cker}, que es un alias de TorrentLocker.
  \item \textbf{Contaminación entre linajes.} Dos notas atribuidas a HELLOKITTY compartían el párrafo de cierre con doce notas de DHARMA y ninguna de su familia declarada. Su retiro eliminó 146 confusiones desde DHARMA y 20 desde PHOBOS en la matriz de confusión.
  \item \textbf{Campaña frente a familia.} Los sucesores de WASTEDLOCKER atribuidos al mismo actor son \textit{familias distintas}, y sus notas no pertenecen al corpus de WASTEDLOCKER aunque la atribución del actor sea correcta.
\end{itemize}

De ahí la regla adoptada: \textbf{el corpus se etiqueta por familia —entendida como el mismo linaje de código y de plantilla— y no por actor ni por campaña}, y ante nombres parecidos se exige una fuente que distinga explícitamente entre ambos. El emparejamiento automático de nombres se evitó en favor de una tabla de equivalencias escrita a mano, porque un normalizador de cadenas es exactamente el mecanismo por el que se cuelan estos tres errores.

\subsubsection*{Fuentes y trazabilidad}

Cada nota lleva registrada su procedencia en un manifiesto: la fuente, la URL cuando existe, el nombre original del archivo cuando se conserva y la forma en que se obtuvo el texto. Las fuentes se admiten con criterios explícitos —autoría verificable, dominio no susceptible de caducar y ser recomprado, preferencia por el proveedor original sobre agregadores comerciales— y se excluyen los repositorios de muestras ejecutables. Al cierre del corpus de 149 notas, \textbf{la totalidad tiene fuente verificable}, frente a las 47 sin procedencia registrada que había en versiones anteriores.

Cuando una nota se transcribe de una imagen, el texto se registra literalmente, incluidas las erratas de la fuente y las truncaduras que ella misma introduce, porque corregirlas produciría un texto que nunca existió. Los identificadores criptográficos transcritos así se verifican además \textbf{mediante su propia suma de comprobación}: una dirección de Bitcoin o de Monero incorpora dígitos de control, de modo que un solo carácter mal leído invalida la verificación. El procedimiento detectó dos caracteres incorrectos en una transcripción externa que a simple vista parecía correcta.

\subsubsection*{Limitación que se declara}

Ninguno de estos controles convierte al corpus en una muestra representativa. Las notas provienen de repositorios públicos, que sobre-representan familias analizadas por la industria y sub-representan campañas recientes o de bajo perfil; el propio proceso de catalogación elimina metadatos —el nombre original del archivo, en particular— que en un incidente real sí están disponibles. Los resultados de este trabajo deben leerse sobre esa base: \textbf{un corpus público, auditado y trazable, pero no una muestra aleatoria del fenómeno}.
